
\section{Aufbau eines Empfängers}
\label{section:aufbau_empfaenger}
\begin{frame}%STARTCONTENT

\frametitle{Grundsätzlicher Aufbau eines Empfängers}
\begin{columns}
    \begin{column}{0.38\textwidth}
        \begin{figure}
            \DARCimage{\linewidth}{736include}
            \caption{\scriptsize Blockdiagramm eines einfachen Empfängers}
            \label{aufbau_empfaenger_blockdiagramm}
        \end{figure}
    \end{column}
    \begin{column}{0.58\textwidth}
        \small
        \begin{enumerate}[<+->]
            \item \textbf{Antenne:} Nimmt Vielzahl von Funkwellen auf und leitet sie als
            elektrische Schwingungen weiter
            \item \textbf{Bandpassfilter:} Lässt nur gewünschten Frequenzbereich
            durch und sperrt alle anderen ungewünschten Frequenzen
            \item \textbf{HF-Verstärker:} Verstärkt das herausgefilterte Signal
            \item \textbf{Demodulator:} Zurückgewinnung des ursprünglichen
            Signals -- Ergebnis ist das Niederfrequenz-Signal (NF)
            \item \textbf{NF-Verstärker:} Verstärkt das demodulierte Signal
            \item \textbf{Lautsprecher:} Wandelt elektrische Schwingung in Schallwellen um
        \end{enumerate}
    \end{column}
\end{columns}
\end{frame}

\begin{frame}
\only<1>{
\begin{PQuestion}{NF201}{Was stellt folgendes Blockdiagramm dar?}{Sender}
{Tongenerator}
{Relaisfunkstelle}
{Empfänger}
{\centering\DARCimage{0.9\linewidth}{525include}}\end{PQuestion}

}
\only<2>{
\begin{PQuestion}{NF201}{Was stellt folgendes Blockdiagramm dar?}{Sender}
{Tongenerator}
{Relaisfunkstelle}
{\textbf{\textcolor{DARCgreen}{Empfänger}}}
{\centering\DARCimage{0.9\linewidth}{525include}}\end{PQuestion}

}
\end{frame}

\begin{frame}
\frametitle{Empfindlichkeit}
\begin{itemize}
  \item Je nach Aufbau haben Empfänger unterschiedliche Eigenschaften
  \item Wichtige Eigenschaft: \emph{Empfindlichkeit}
  \item Fähigkeit, schwache Signale zu empfangen
  \item Je empfindlicher, umso schwächere Signale können empfangen werden
  \end{itemize}
\end{frame}

\begin{frame}
\only<1>{
\begin{QQuestion}{NF303}{Worauf bezieht sich die Empfindlichkeit eines Empfängers?}{Auf die Fähigkeit, schwache Signale zu empfangen}
{Auf die Stabilität des VFO}
{Auf die Bandbreite des HF-Vorverstärkers}
{Auf die Fähigkeit, starke Signale zu unterdrücken}
\end{QQuestion}

}
\only<2>{
\begin{QQuestion}{NF303}{Worauf bezieht sich die Empfindlichkeit eines Empfängers?}{\textbf{\textcolor{DARCgreen}{Auf die Fähigkeit, schwache Signale zu empfangen}}}
{Auf die Stabilität des VFO}
{Auf die Bandbreite des HF-Vorverstärkers}
{Auf die Fähigkeit, starke Signale zu unterdrücken}
\end{QQuestion}

}
\end{frame}%ENDCONTENT
